\section{结论}
\label{sec:conclusion}

RoPE 并不会对远程证据施加简单的单调惩罚。它使固定的内容匹配变为距离相关相位的
混合；不利的干涉会在 softmax 之前降低其 QK 优先级。随之改变的注意力写入又会改变
后续的 pre-RoPE 表征，并可能跨越输出决策边界。精确的首层重建、逐层有限差分重放和
激活修补共同支持 Qwen3-8B 中的这一机制。同一诊断也导出了一种保守的修复方案：保留
局部 RoPE 和精确的 post-RoPE 候选使用，同时为远程 token 加入无位置信息的语义候选提议。
\methodpost 在留出的受控两跳检索任务上取得了改进，其中在 16K--32K 长度下最为明显。
要确立其普适性，仍需在自然基准、更多架构、局部顺序控制和高效检索索引上开展验证；
这些均是当前论断的明确边界。
